% INTRODUCCIÓN

\section{Introducción}

La calidad del aire se ha consolidado como un problema técnico, sanitario y de gestión pública que exige sistemas de monitoreo capaces de producir información útil para decisiones operativas en tiempo real. En este TFM, ese reto se aborda mediante la propuesta titulada \PortadaTitulo, que mide contaminantes y variables ambientales, consolida el AQI con base normativa EPA/AQS, aplica una inferencia difusa principal para interpretar el riesgo y complementa esa salida con una capa contextual basada en reglas para modularla cuando existen variables ambientales disponibles. En este contexto, los índices de calidad del aire (AQI) han sido ampliamente utilizados para traducir concentraciones de contaminantes en categorías interpretables por autoridades y ciudadanía. Sin embargo, revisiones de amplio alcance reportan una fragmentación metodológica persistente: diferencias en contaminantes considerados, escalas, umbrales, funciones de agregación y criterios de interpretación entre países y ciudades, lo que reduce la comparabilidad y genera ambigüedad en la lectura del riesgo \cite{ART01,ART02}. En particular, se ha documentado que varios enfoques subestiman condiciones reales cuando priorizan un único contaminante dominante, y que muchas formulaciones no incorporan adecuadamente efectos combinados entre contaminantes ni escenarios de exposición de baja intensidad pero sostenida \cite{ART01}.

Sobre esa base, la literatura reciente muestra un desplazamiento hacia enfoques de lógica difusa para tratar incertidumbre, imprecisión y no linealidad en datos ambientales. Esta línea tiene sustento conceptual en el marco de variables lingüísticas y reglas difusas para sistemas complejos y difícilmente modelables con esquemas estrictamente deterministas \cite{ART03}. En aplicaciones contemporáneas de monitoreo de aire, se observan resultados prometedores en clasificación, soporte a la decisión y, en algunos casos, predicción usando FIS, ANFIS y variantes híbridas \cite{ART04,ART05,ART06,ART07,ART08}. No obstante, la evidencia también revela divergencias relevantes entre autores: algunos trabajos privilegian precisión de clasificación en datasets concretos \cite{ART04}, otros priorizan control ambiental en espacios cerrados y automatización de ventilación \cite{ART05}, otros enfatizan implementación IoT con desempeño operativo local \cite{ART06,ART07}, mientras que propuestas más recientes se orientan a eficiencia de red y enrutamiento inteligente en WSN para ciudades inteligentes \cite{ART08}. Esta heterogeneidad confirma que el campo está activo, pero todavía disperso en sus criterios de diseño y validación.

En relación con los vacíos y limitaciones explícitamente reportados por los propios estudios, persisten cuatro problemas de fondo. Primero, la falta de estandarización y trazabilidad metodológica para comparar enfoques de forma justa entre contextos, sensores y configuraciones experimentales \cite{ART01,ART02}. Segundo, limitaciones de generalización: varios desarrollos se validan en entornos localizados o con cobertura restringida, con necesidad de pruebas en escenarios más diversos y escalables \cite{ART04,ART08}. Tercero, restricciones técnicas que afectan despliegue real, como calibración de sensores de bajo costo, deriva de medición, latencia, consumo energético, y ausencia de mecanismos robustos de seguridad en arquitectura distribuida \cite{ART08}. Cuarto, una brecha entre el desempeño reportado en simulación/prototipo y la documentación reproducible de decisiones de diseño, criterios de evaluación y robustez estadística \cite{ART06,ART07,ART08}. Estos vacíos dificultan pasar de soluciones puntuales a plataformas confiables y transferibles.

Desde esta perspectiva, la motivación de esta investigación es doble. Por un lado, se requiere una síntesis rigurosa de evidencia reciente que permita identificar con precisión qué variables, diseños difusos y arquitecturas de plataforma han mostrado mayor utilidad práctica en monitoreo de calidad del aire en tiempo real. Por otro lado, esa síntesis debe traducirse en decisiones de ingeniería para una propuesta técnica concreta, evitando que el estado del arte quede desconectado de la implementación. En consecuencia, el trabajo describe antecedentes y busca convertir evidencia dispersa en criterios de diseño trazables para una integrarlos en una solución aplicable.

A partir de los antecedentes y de los vacíos observados, este estudio establece como brecha central la ausencia de un marco integrador que conecte, de manera explícita y reproducible, cuatro niveles que suelen tratarse por separado: (i) selección y tratamiento de variables ambientales y contaminantes, (ii) diseño del sistema de inferencia difusa y su rol en la toma de decisiones, (iii) evaluación flexible del riesgo mediante términos lingüísticos, y (iv) arquitectura de monitoreo en tiempo real con capacidad de alertamiento preventivo bajo condiciones reales de despliegue. Esta brecha es crítica porque afecta a la calidad de la revisión sistemática y a la viabilidad técnica de la plataforma que se propone desarrollar.

Con ese marco, el objetivo general que guía todo el trabajo es desarrollar un módulo de monitoreo de calidad del aire para apoyar la evaluación del riesgo y la generación de alertas, mediante inferencia difusa y datos abiertos de redes de monitoreo ambiental, en entornos urbanos. En coherencia con ello, la pregunta general de investigación del TFM es: ¿cómo desarrollar un módulo de monitoreo de calidad del aire que apoye la evaluación del riesgo y la generación de alertas mediante inferencia difusa y datos abiertos de redes de monitoreo ambiental en entornos urbanos? La fase de revisión sistemática formula, además, una pregunta específica orientada a caracterizar el estado del arte y a derivar criterios de diseño para la propuesta.

Bajo esta orientación, las investigaciones referentes se interpretan como complementarios más que excluyentes. Los estudios de revisión de AQI aportan diagnóstico de inconsistencias estructurales y criterios para evitar comparaciones engañosas \cite{ART01,ART02}; los trabajos de base difusa aportan fundamento para modelar incertidumbre en dominios complejos \cite{ART03}; y los desarrollos aplicados en IoT/WSN muestran rutas tecnológicas concretas para captura, inferencia y operación en campo \cite{ART04,ART05,ART06,ART07,ART08}. La contribución de este TFM es precisamente cerrar el ciclo entre esos tres planos: evidencia, método y propuesta técnica.

El documento se estructura de la siguiente manera. Tras esta introducción, la sección de objetivos delimita el alcance del estudio y los resultados esperados. Luego, en el estado del arte y marco teórico se presenta la síntesis derivada de la revisión sistemática. A continuación, en desarrollo del proyecto y resultados se describe la metodología aplicada, incluyendo la revisión sistemática como primera fase del trabajo, junto con el diseño de la solución, su implementación y el esquema de validación. Finalmente, la sección de conclusiones y trabajos futuros recoge los principales hallazgos, limitaciones del estudio y líneas de continuidad para la evolución de la propuesta.
